\definition{Etimologías}{
    \begin{itemize}
        \item Grasa: Del latín \textit{crassus}, que significa grueso, denso o sucio.
        \item Lípido: Del griego \textit{lipos}, que se refiere a una grasa para alimento o grasa para unciones sagradas.
        \item Aceite: Del latín \textit{oleo}, que a su vez se deriva del griego \textit{elaka} y \text{olivo}, se refiere al árbol del que se obtiene el \textit{rey de los aceites}.
    \end{itemize}
}

\section{Estructura}

Característica principal: Solubles en solventes orgánicos

\begin{figure}[H]
    \centering
    \chemfig{R-[:30]C-[:-30]O-[:30]C=[:-30]O}
    \caption{Estructura general de un lípido}
    \label{fig:estructura_lípidos}
\end{figure}

\begin{itemize}
    \item No forman polímeros
    \item Son pequeñas
    \item Se asocian por fuerzas covalentes
    \item Cuentan con una parte hidrofílica y una hidrofóbica
\end{itemize}

\section{Micelas}

A partir de su característica estructural de contar con una parte hidrofílica y una hidrofóbica, los lípidos pueden formar estructuras circulares llamadas micelas, las cuales, a través de fuerzas polares, fuerzas no polares y fuerzas de Van der Waals, se estructuran colocando la parte hidrofílica de los lípidos hacia afuera y la parte hidrofóbica hacia dentro. Estas micelas son de utilidad en los productos de limpieza.

\section{Funciones}

Los lipidos, por sus características estructurales, cuentan con varias funciones:

\begin{itemize}
    \item Hormonal: estrógeno, vitaminas (A, D, K y E)
    \item Estructural: Membrana celular
    \item Energía: Tejido adiposo
    \item Protección térmica
\end{itemize}